% Source: https://artofproblemsolving.com/community/c3h3166099_2019_lmt_spring_guts_round_p25p36__lexington_mathematical_tournament
% Stable identity: aops-topic-3166099-q33
Put $b=f(0)$, $c=g(0)$, $F(x)=f(x)-b$, and $G(y)=g(y)-c$. Comparing the equation with its case $y=0$ gives $F(x+G(y))=F(x)+y$. In particular $F(G(y))=y$, so $G$ is injective and $q=G(1)\ne0$. For any rational $y$, write $\dfrac{G(y)}{q}=\dfrac{m}{n}$ with integers $m$ and $n>0$. Repeated translation by $G(y)$ changes $F$ by $ny$, whereas translation by $q$ repeated $m$ times changes it by $m$; these are the same shift, so $ny=m$. Thus $G(y)=qy$ for every rational $y$, and surjectivity of these shifts gives $F(x)=\dfrac{x}{q}$. Write $a=\dfrac{1}{q}\in\mathbb Q\setminus\{0\}$. Then $f(x)=ax+b$, $g(y)=\dfrac{y}{a}+c$. The original equation at $y=0$ forces $h(t)=at+ac(a-1)$, because $f$ is surjective. The first two supplied values give $b=2-6a$ and $c=2-\dfrac{1}{2a}$. The third becomes $\dfrac{7a}{2}+(2a-\dfrac{1}{2})(a-1)=2$, or $4a^2+2a-3=0$. Its roots $\dfrac{-1\pm\sqrt{13}}{4}$ are irrational, contradicting $a\in\mathbb Q$. Therefore no such functions exist and the set of possible values is $\boxed{\varnothing}$.
